% ---------------------------------------------------------------------
%  Annexes (10 pages MAX — PAS de code complet, cf. consignes)
%  Le code source complet n'est PAS reproduit ici : il est référencé
%  sur une page web (cf. \ref au dépôt dans le corps du texte).
% ---------------------------------------------------------------------
\chapter{Annexes}
\label{chap:annexes}

\section{Extraits de commandes (lignes principales)}
\label{ann:code}
% Les consignes autorisent les « lignes principales des commandes »,
% jamais le programme complet. Exemples représentatifs ci-dessous.

\begin{lstlisting}[caption={Test de l'avantage moyen (extrait).}]
# Shapiro (normalité), Student 1-échantillon, seuil Bonferroni
from scipy import stats
w, p_shapiro = stats.shapiro(edge.sample(min(len(edge), 5000)))
t, p_student = stats.ttest_1samp(edge, 0, alternative="greater")
seuil = 0.05 / n_strategies      # Bonferroni
\end{lstlisting}

\begin{lstlisting}[caption={ANOVA à deux facteurs avec interaction (extrait).}]
import statsmodels.api as sm
from statsmodels.formula.api import ols
modele = ols("edge ~ C(strategy) * C(regime)", data=df).fit()
sm.stats.anova_lm(modele, typ=2)
\end{lstlisting}

\section{Tables complémentaires}
\label{ann:tables}

\subsection*{Test de normalité de Shapiro--Wilk (étape 1)}

Le tableau~\ref{tab:ann-shapiro} détaille, pour chaque stratégie, la
statistique \(W\) et la p-value du test de Shapiro--Wilk appliqué à
l'edge. La normalité est rejetée pour les dix stratégies. On remarque que
même les stratégies dont le \(W\) est proche de 1 (par exemple
\texttt{dt\_top}, \(W = 0{,}95\)) sont rejetées~: à grand \(n\), le
moindre écart à la normale suffit à produire une p-value négligeable, ce
qui justifie le recours au théorème central limite plutôt qu'à ce test
pour valider l'emploi de Student.

\begin{table}[H]
\centering
\caption{Statistique \(W\) et p-value de Shapiro--Wilk par stratégie
(edge~; sous-échantillon de 5\,000 trades lorsque \(n > 5\,000\)).}
\label{tab:ann-shapiro}
\small
\begin{tabular}{@{}lrrr@{}}
\toprule
Stratégie & $n$ & $W$ & $p$ (Shapiro) \\
\midrule
\texttt{db\_bottom}    & 8\,513 & $0{,}911$ & $2{,}7\times10^{-47}$ \\
\texttt{dt\_top}       & 7\,381 & $0{,}950$ & $3{,}2\times10^{-38}$ \\
\texttt{hs\_classic}   & 1\,065 & $0{,}886$ & $2{,}4\times10^{-27}$ \\
\texttt{hs\_inverse}   & 1\,206 & $0{,}963$ & $6{,}4\times10^{-17}$ \\
\texttt{ma\_crossover} & 5\,417 & $0{,}498$ & $3{,}2\times10^{-80}$ \\
\texttt{rsi\_classic}  & 7\,117 & $0{,}864$ & $1{,}6\times10^{-54}$ \\
\texttt{rsi\_strict}   &   978  & $0{,}717$ & $1{,}1\times10^{-37}$ \\
\texttt{rsi\_trend}    & 2\,789 & $0{,}883$ & $1{,}2\times10^{-41}$ \\
\texttt{sr\_breakdown} & 1\,377 & $0{,}775$ & $8{,}0\times10^{-40}$ \\
\texttt{sr\_breakout}  & 2\,367 & $0{,}843$ & $2{,}1\times10^{-43}$ \\
\bottomrule
\end{tabular}
\end{table}

\emph{[Autres tableaux détaillés (paires de Tukey, coordonnées
factorielles, etc.) à ajouter au fil des étapes, dans la limite de 10
pages.]}
